\section{Related Work}
\label{sec:related_work}
% Robotic simulation is pivotal for scaling up robot demonstrations \cite{o2024open}, enabling robust policy evaluation \cite{li2024evaluating,team2025evaluating}, and accelerating reinforcement learning \cite{ebert2018visual} without the risks and costs of real-world deployment. 
We review the evolution of embodied simulation, tracing its trajectory from classical physics engines to the recent emergence of generative world models.

\paragraph{Physical simulation.}
Classical robotic simulation methods rely on physics engines, such as MuJoCo, IssacSim, SAPIEN \cite{todorov2012mujoco,erez2015simulation,xiang2020sapien,makoviychuk2021isaac,mittal2023orbit,robotwin}. They simulate interactions based on rigid-body dynamics, requiring meticulously hand-crafted meshes, precise physical properties (\eg, friction, mass), and pre-defined rules. 
% While they provide high physical fidelity, they often suffer from a “sim-to-real” gap due to the lack of visual realism. To alleviate this, 
Recent real-to-sim advancements further leverage 3D Gaussian Splatting (3DGS) with continuum mechanics or rigid-body solvers to reconstruct digital twins directly from real-world captures \cite{pashevich2019learning,li2024robogsim,abouchakra-embodiedgaussians,jiang2025gsworld,han2025re}. 
% These methods integrate 3DGS with continuum mechanics or rigid-body solvers to enable photo-realistic and physically-driven simulations. 
A few methods also focus on simulating realistic depth \cite{xu2025stable,liu2025manipulation}.
Nevertheless, they rely on explicit physical solvers and predefined laws, which struggle to generalize to complex environmental responses, limiting their scalability across unstructured environments.

\paragraph{Learning world models.}
World Models \cite{allen1983planning,ha2018world} aim to internalize the underlying dynamics of the environment, allowing agents to “dream” and plan within a learned latent space. With the advent of diffusion models \cite{song2019generative,ho2020denoising,dit} and large-scale video pre-training \cite{hong2022cogvideo,bruce2024genie,zhang2025show,wan2025wan}, video generative models have demonstrated an extraordinary ability to capture complex visual dynamics \cite{assran2025v,chefer2025videojam}. Furthermore, \textit{interactive} video generation has advanced through the integration of various conditioning mechanisms \cite{zhang2023adding} designed to represent diverse control signals. Early efforts in this direction have focused on injecting viewpoint trajectories \cite{yu2025wonderworld,liang2025wonderland,ren2025gen3c,gu2025diffusion,Yu_2025_ICCV} to synthesize immersive, navigable environments.

\paragraph{Embodied video-generation models.}
Recently, research has pivoted toward using embodied actions as controls for interactive video generation, aiming to synthesize robot-world interactions with both visual and physical fidelity. Based on the representation of embodied actions at the input level, this research stream can be categorized into four sub-groups.
1) Text: A main-stream works \cite{robodreamer,yang2024learning,robomaster,chen2025world4omni,yang2025roboenvision} generate visual outcomes from high-level text instructions, which lack the fine-grained precision required for low-level manipulations.
2) Latent embedding: Another type of works \cite{wu2024ivideogpt,cen2025worldvla,cen2025rynnvla,huang2025ladi,luo2025learning,zhang2025imowm,zhu2025irasim,agarwal2025cosmos,gao2026dreamdojo,liao2026genie,guo2026ctrl} such as IRASim \cite{zhu2025irasim} and Ctrl-World \cite{guo2026ctrl} typically encode 7-DoF end-effector poses into compressed embeddings. This forces the generative model to infer or `guess' the underlying robot kinematics, often resulting in physically implausible failures.
3) Semantics: ORV \cite{yang2026orv} introduces 3D semantic occupancy of static environments; however, it needs an off-the-shelf occupancy-prediction model to acquire and lacks temporal dynamics, still requiring action encodings or text to provide dynamic information.
4) 2D visual prompt: EVAC \cite{jiang2025enerverse}, AnchorDream \cite{ye2025anchordream}, and VAP \cite{VAP_2025_ICCV} represent embodied actions using 2D prompts like angle arrows, 2D renders, and skeletons. A concurrent work, BridgeV2W \cite{chen2026bridgev2wbridgingvideogeneration}, utilizes URDF to drive robots but merely renders the resulting trajectory into 2D binary masks. Such 2D-based signals lack spatiotemporal constraints and struggle to provide precise control guidance.
Commonly, as for the output, all aforementioned methods use video generations to predict 2D (\ie, RGB) frames, treating the robot/environment as a monolithic pixel stream. Due to the lack of spatiotemporal awareness, they struggle to simulate complex physical interactions. 
%—such as delicate manipulations or material deformations. 

\textbf{Summary}: Whether evaluated by action conditioning or output modality, the aforementioned methods fail to resolve the \textit{\textbf{challenging trilemma}} of \textit{dynamics, precision} and \textit{spatiotemporal awareness}—the three pillars of reliable embodied simulation. To overcome this, our Kinema4D learns intricate dynamics through a 4D generative model, where abstract action controls are grounded via kinematics, securing precision and spatiotemporal awareness, which in turn anchors the generative model to synthesize complex dynamics.
% bypassing the limitations of explicit physical solvers while exceeding the spatial consistency of 2D video models.
% To overcome this, we propose 4D-native grounding. We employ kinematics to transform abstract robot actions into pointmap sequences, enjoying a deterministic link between 3D physical reality and generative synthesis, enabling the precise modeling of complex, reactive dynamics.
% there is no inherent mechanism to ensure that the synthesized environment remains consistent with the robot’s 4D pose.

\paragraph{3D/4D world models.}
Several recent works have been able to output 3D/4D worlds. First, ParticleFormer \cite{huang2025particleformer} and PointWorld \cite{huang2026pointworld} directly utilize simple transformers to predict the trajectories of pre-defined 3D particles for both robots and environments. However, they lack the generative flexibility required to synthesize new geometry that depicts emergent world dynamics beyond the initial 3D input.
On the other hand, 4D-native generative models that aim to generate full spatio-temporal sequences have gained significant attention. For instance, Aether \cite{Aether} unifies geometry-aware reasoning by jointly optimizing 4D dynamic reconstruction and goal-conditioned video prediction, while 4DNex \cite{chen20254dnex} enables 4D world generation in a single feed-forward pass. 
% We extend the potent capabilities of these 4D-native architectures to the domain of robotic simulation.
Although recent advancements such as TesserAct \cite{zhen2025tesseract}, GWM \cite{lu2025gwm}, iMoWM \cite{zhang2025imowm}, and Robo4DGen \cite{robo4dgen} demonstrate the ability to synthesize 4D embodied worlds at the output stage, they remain constrained by a critical design bottleneck: they solely inject \textit{static} guidance (\eg, depth maps, surface normals, or Gaussian splats) from the initial world environment, which \textit{lacks temporal dynamics}. Consequently, these models still rely on text instructions or latent tokens to inject robot actions, which \textit{lack the granularity} required for fine-grained interactions. In contrast, by transforming abstract robot actions into 4D pointmaps, our Kinema4D accepts \textit{precise} robot controls with \textit{both spatial and temporal cues}. To fully utilize such controls, we apply a 4D-native generative model with spatiotemporal reasoning abilities, allowing the model to not only understand robot movement but also to flexibly predict its physical effects on the environment for high-fidelity simulation.
% By predicting full 4D sequences, our framework maintains geometric rigor throughout the entire generative process, 
% Unlike heuristic 2D masks that offer only coarse spatial hints, pointmaps provide a pixel-aligned, depth-aware representation of 3D occupancy.